% DASFAA/LNCS camera-ready template for JointTune
\documentclass[runningheads]{llncs}

\usepackage[T1]{fontenc}
\usepackage{newtxtext}
\usepackage[varvw]{newtxmath}
\usepackage{graphicx}
\usepackage{amsmath}
\usepackage{amssymb}
\usepackage{booktabs}
\usepackage{algorithm}
\usepackage{algpseudocode}
\usepackage{url}

\begin{document}

\title{JointTune: Jointly Revising Search Structure and Candidate
Allocation for LLM-Assisted Database Tuning}

\titlerunning{JointTune: Joint Search-State Revision and Allocation}

\author{RenJie Yuan\inst{1}\orcidID{0009-0008-7164-7879}}

\authorrunning{R. Yuan}

\institute{University of Electronic Science and Technology of China, Chengdu, China\\
\email{yuanrenjie@example.com}}

\maketitle

\begin{abstract}
Database knob tuning under a limited execution budget must jointly
decide which parameters to tune, where to explore their values, and how
much candidate-generation capacity each direction receives. Large
language models (LLMs) can supply parameter semantics, recommended
regions, and interaction hypotheses, but these priors must be corrected
by execution evidence for the target workload and hardware. Existing
LLM-assisted tuners use execution feedback mainly to update performance
prediction, while the subsequent search structure and the share of
expert proposals are revised only partially or independently. This paper
presents JointTune, a framework that couples search-state revision with
audited candidate allocation in a single feedback loop. Stage~1 compares
a bounded menu of keep, reshape, exchange, and combined actions, where
each action induces a complete search state: its regions determine audit
points, audit opportunities determine quotas, and these quotas enter the
action score. Experts review the resulting posterior evidence and return
legal patches. Stage~2 allocates real execution budget on the selected
state through equal-size audits and informative trial selection, where a
trial may be chosen for its value in clarifying which direction deserves
the next quota. Successful executions update the shared Gaussian process,
local coverage, residuals, and audit credit, whereas failed attempts
remain in history without training the utility GP. The three search
decisions share the same execution evidence and posterior model, so
limited evaluations revise both what to search and how much to invest,
rather than optimizing them separately and then combining them. We
present the method design, its key invariants, and an experimental
protocol; systematic empirical results will be reported in the completed
study.
\keywords{Database tuning \and Large language models \and Joint
search-state revision \and Audited allocation \and Bayesian optimization
\and Candidate budget allocation.}
\end{abstract}

\section{Introduction}

The growth of data-intensive applications and increasingly diverse
workloads places greater demands on database performance and resource
efficiency. Database knob tuning is a critical means of fully utilizing
database resources: by adjusting configuration parameters governing
memory, concurrency, storage, and query optimization, it improves
throughput or latency for a given workload and environment. However,
modern database management systems (DBMSs) typically expose hundreds of
heterogeneous and interdependent knobs, forming a vast configuration
space. Since evaluating a configuration requires executing the target
workload, extensive trial-and-error search incurs high time and resource
costs. Therefore, the core problem of efficient tuning is: \textbf{how to
construct and continuously adjust an effective search space focused on
promising directions under a limited execution budget.}

Existing approaches broadly fall into three categories: heuristic,
machine-learning-based, and large language model (LLM)-based methods.
\textbf{Heuristic methods} use predefined rules or search strategies, as
exemplified by BestConfig~\cite{bestconfig} and
OpenTuner~\cite{opentuner}, but their effectiveness depends heavily on
prior assumptions about the tuning task. \textbf{Machine-learning-based
methods} use Bayesian optimization or reinforcement learning to learn
from execution feedback~\cite{ottertune,ituned,cdbtune,qtune}; however,
exploring high-dimensional spaces usually requires substantial evaluation
effort. \textbf{LLM-based methods} incorporate domain knowledge into knob
selection, range reduction, and configuration recommendation to narrow
the search space and accelerate
tuning~\cite{gptuner,mctuner,agenttune}. SysInsight~\cite{sysinsight}
further explores extracting verifiable knowledge from database source
code and runtime diagnostics to guide tuning. These approaches
demonstrate that domain knowledge can effectively limit the search space
before tuning begins. However, documentation often lacks explicit tuning
advice for every knob, and most recommendations are broad and may not
suit specific workloads, hardware conditions, and configuration
combinations.

Some studies have also begun to use execution feedback to adjust search
arrangements. DOT~\cite{dot} updates the active knob set during tuning
through statistical tests, and MCTuner~\cite{mctuner} organizes local
search through spatial decomposition. However, these methods typically
revise only one aspect of the search space. DOT's decisions mainly act on
the inclusion and removal of knobs, while value regions and candidate
budgets are handled implicitly by the downstream optimizer; when low
gains arise from a misaligned region or an insufficient budget,
adjusting the knob set alone cannot recover search efficiency. MCTuner
incorporates LLM knowledge for knob selection and range reduction, but
its region revision and candidate allocation still follow a predefined
spatial decomposition rather than execution-driven joint revision. These
methods therefore treat search-space adjustments as independent aspects,
ignoring the inseparable coupling among parameter groups, recommendation
regions, and candidate-generation budgets. Introducing new parameters
without adjusting their value regions may leave them wandering in
ineffective areas; expanding regions without increasing candidate budgets
leads only to ineffective exploration; and increasing budgets without
updating parameters and regions only directs resources toward
low-potential directions. Without a unified feedback loop, these isolated
modifications cannot effectively translate into an improvement of the
actual search space.

To address this, we formulate the tuning process as follows: under a
limited budget, each round jointly determines the active parameter set
\(A_t\), the recommendation region \(R_t\), and the candidate generation
budget \(b_t\) based on execution evidence, aiming to maximize performance
improvement under finite evaluations. These three elements collectively
determine which configurations can be proposed and evaluated. Around this
joint revision goal, we face three interrelated and coupled challenges:

\textbf{C1: Configuration-dependent parameter value identification and
its coordination with regions and budgets.} Whether a parameter is worth
continuing to search depends not only on its own potential but is also
constrained by whether its current recommendation region is misaligned
and whether sufficient candidate budget supports it. Low local gains may
stem from the parameter being exhausted, or from region misplacement or
budget insufficiency. The system must jointly determine parameter
retention and its supporting budget under limited observations. Unlike
dynamic-selection methods that adjust only the active knob
set~\cite{dot}, our parameter-retention decision also accounts for the
region state and budget support, preventing a low gain caused by region
misplacement from being misread as parameter exhaustion.

\textbf{C2: Matching prior recommendation regions with actual
distributions under budget constraints.} Documentation and LLMs often
lack explicit tuning advice for a specific knob, and most recommendations
are broad and may not suit the current workload, hardware, and
configuration combination. Effective region revisions, such as shifting
or expanding, require corresponding candidate-budget support; blindly
expanding a region without allocating sufficient candidates leads only to
ineffective exploration. The system must balance region refinement,
expansion, and budget allocation based on execution evidence.

\textbf{C3: Cross-domain knowledge coverage and fair allocation of
limited budgets.} Database tuning spans multiple specialized domains, and
a single agent can hardly cover the configuration rules of all domains,
making biased or missed judgments likely. Introducing multiple domain
experts mitigates this, but their suggestions may overlap or conflict. If
evaluation budgets are allocated by generation volume, some valuable
parameter directions may be misjudged for lack of evidence, disrupting
the feedback loop between parameter and region revisions. The system must
ensure fair evaluation budgets so that feedback evidence can effectively
guide the next round of joint revision.

To solve these challenges, we propose \textbf{JointTune}, an LLM-assisted
database tuning framework that jointly revises the effective search space
through execution feedback. JointTune organizes multi-domain experts to
propose configuration suggestions and uses real execution evidence to
uniformly model the active parameters, recommendation regions, and
candidate budgets as a search-space state updated with evidence, enabling
the domain-knowledge-guided search to gradually adapt to the target
environment.

Our main contributions are as follows:

\begin{enumerate}
    \item \textbf{An execution-feedback-driven joint search-space revision
    framework.} We uniformly model the active parameter set,
    recommendation regions, and candidate generation budgets as a
    search-space state updated with evidence. Within this framework, we
    design mutually coupled mechanisms to address the decoupling between
    the search space and budget caused by independently revising any
    single element.

    \item \textbf{A conditional-potential-driven dynamic parameter
    activation and region reshaping mechanism.} This mechanism combines
    the potential gains under the current configuration with actual
    exploration evidence to dynamically adjust the active parameter set.
    It revises recommendation regions based on execution feedback,
    enabling the search to flexibly switch between local refinement and
    cross-region exploration, breaking the limitations of static prior
    ranges.

    \item \textbf{A feedback-guided multi-expert equal-size budget
    allocation mechanism.} This mechanism utilizes domain experts to
    generate complementary configurations and employs equal-size audits
    to reduce the impact of candidate-quantity differences on opportunity
    comparison. The updated budgets are accurately fed back into the
    revision of parameters and regions, ensuring that feedback evidence
    can fairly guide the next round of joint optimization.
\end{enumerate}

\section{Related Work}

\subsection{Heuristic Methods}
Early automated tuning systems predominantly relied on heuristic rules
or predefined search strategies. Rule-based tools, such as
PGTune~\cite{pgtune}, directly derive configurations based on hardware
specifications. Search-based methods like OpenTuner~\cite{opentuner}
integrate multiple search algorithms, while BestConfig~\cite{bestconfig}
combines sampling with recursive divide-and-conquer. Although these
methods are lightweight and easy to deploy, their effectiveness heavily
depends on manually designed rules or strong assumptions about the tuning
task. Consequently, they struggle to capture the highly non-linear
interactions among knobs under complex workload conditions.

\subsection{Bayesian Optimization-Based Methods}
Bayesian Optimization (BO)-based methods approximate the performance
landscape using surrogate models to balance exploration and exploitation.
Specifically, iTuned~\cite{ituned} applies Gaussian Processes,
OtterTune~\cite{ottertune} integrates factor analysis and transfer
learning, while SMAC~\cite{smac} and ResTune~\cite{restune} adopt Random
Forests and meta-learning, respectively. To address high-dimensional
spaces, LlamaTune~\cite{llamatune} employs random projections and knob
bucketing. Recently, BO has been extended to dynamic scenarios:
OnlineTune~\cite{onlinetune} introduces contextual BO with spatial
partitioning, and SafeTune~\cite{safetune} filters high-risk
configurations via semi-supervised anomaly detection. Notably,
DOT~\cite{dot} combines Recursive Feature Elimination with
Cross-Validation (RFECV) and a likelihood ratio test to dynamically prune
and expand the active knob set during tuning, representing the
dynamic-selection route. However, DOT's decisions mainly act on knob
inclusion and removal, while value regions and candidate budgets are
handled implicitly by the downstream optimizer rather than being treated
as explicit joint decision variables; when the effective region shifts
across configurations or the candidate budget is insufficient to cover
that region, statistical tests on the knob set alone cannot distinguish
an exhausted parameter from a misaligned region or an insufficient
budget. This differs from our goal of jointly revising active knobs,
value regions, and candidate budgets. Despite these advances, building
reliable surrogate models in high-dimensional spaces typically still
requires hundreds of real workload evaluations.

\subsection{Reinforcement Learning-Based Methods}
Reinforcement Learning (RL) formulates tuning as a Markov Decision
Process to learn configuration policies via interactions.
CDBTune~\cite{cdbtune} and QTune~\cite{qtune} utilize Deep Deterministic
Policy Gradient (DDPG), with QTune further incorporating query features
to enhance decision-making. HUNTER~\cite{hunter} and
WATuning~\cite{watuning} improve policy generalization through rule
extraction and attention mechanisms, respectively. To handle dynamic
workloads, MR-MDRL~\cite{mrmdrl} constructs reward signals from native
runtime metrics for dynamic multi-objective trade-offs, while
CTuner~\cite{ctuner} proposes a causal RL framework for NoSQL databases.
Nevertheless, RL methods still suffer from training instability and low
sample efficiency in large-scale knob spaces.

\subsection{LLM-Based Database Tuning Methods}
Recently, Large Language Models (LLMs) have been widely introduced into
database tuning to alleviate cold-start issues and high search costs.
Many studies utilize LLMs to extract and organize structured knowledge to
reduce the search space. For example, GPTuner~\cite{gptuner} and
DB-BERT~\cite{dbbert} extract knowledge from manuals and forums to filter
redundant knobs and recommend values. MCTuner~\cite{mctuner} employs a
Mixture-of-Experts (MoE) mechanism to evaluate key knobs and reduce
ranges, subsequently delegating the search to hierarchical spatial
decomposition and BO. KnobCF~\cite{knobcf} and Rabbit~\cite{rabbit}
organize reliable knowledge through uncertainty-aware classification and
retrieval-augmented generation (RAG), respectively, to minimize
ineffective evaluations.

Furthermore, some works enable LLMs to participate more deeply in
candidate generation and search control. AgentTune~\cite{agenttune}
decomposes the tuning task into multiple agents and uses tree search to
directly generate candidate configurations. E2ETune~\cite{e2etune}
fine-tunes a generative model to predict configurations end-to-end based
on workload characteristics. Additionally, recent studies like
SysInsight~\cite{sysinsight} explore extracting verifiable tuning clues
from database source code and runtime diagnostics to guide the parameter
tuning process.
\section{Problem Definition}\label{sec:problem}

Let \(K\) be the set of tunable knobs in a database management system (DBMS), and a configuration vector be \(x = (x_1, \ldots, x_d)\), where \(d=|K|\). All configurations satisfying database resource limits, type constraints, and activation conditions form the legal search space \(\mathcal{X}\).

Given a specific workload and hardware environment, evaluating a configuration \(x\) requires executing the target workload, and its true performance metric is represented by a function \(f(x)\). To unify the optimization direction, we transform throughput or latency into a utility \(u(x)\) to be maximized (e.g., transforming the latency into a normalized negative value, or converting throughput into a logarithmic value).

The goal of database parameter tuning is to find the optimal configuration that maximizes the utility under a limited real execution budget \(T\):
\begin{equation}
x^* = \arg\max_{x \in \mathcal{X}} u(x), \quad \text{s.t.} \quad \text{Cost}(x) \le T.
\end{equation}

\section{Method}\label{sec:method}

\subsection{Overview}\label{sec:overview}

The overall workflow of JointTune is illustrated in
Figure~\ref{fig:ct-framework} and comprises two core stages.

\noindent\textbf{Stage 1: Evidence-guided search-state revision.}
The first stage addresses how the search structure should change given
limited execution evidence. JointTune fits a shared Gaussian process (GP)
on successful observations and forms a bounded action menu: keep the
current structure, reshape an incumbent region, exchange an eligible group
for a stronger inactive one, or combine exchange with reshaping a retained
partner. Each action induces a complete search state: its regions determine
audit points, audit opportunities determine quotas, and these quotas enter
the action score. Experts review the posterior evidence in their regular
requests and return legal patches and region suggestions. The controller
builds equal-size audits for every candidate state and selects a state only
when the predictive model is calibrated and audits are available; otherwise
it retains the current structure. The output of this stage is a revised
active knob set, value regions, and mode quotas.

\noindent\textbf{Stage 2: Audited allocation and informative trial selection.}
The second stage addresses where limited candidate capacity should be
invested on the revised structure. The selected state generates candidates
through group modes and reserved reference, cross-domain, and coverage
channels. The controller deduplicates the candidates, scores them with
noisy expected improvement, and may replace the highest-EI candidate by one
with higher information about which direction deserves the next quota;
scheduled coverage and re-entry probes override this replacement. After one
candidate is executed, a successful outcome updates the GP, local coverage,
residuals, and audit credit, whereas a failed attempt remains in history but
does not train the utility GP. The updated evidence enters the next round of
state revision, closing the feedback loop.

JointTune couples search-space revision with candidate-capacity allocation
in a single feedback loop: Stage~1 uses posterior evidence to compare
retaining, reshaping, and replacing, deciding what to search and where;
Stage~2 allocates real execution budget on the selected structure through
equal-size audits and informative trials, deciding how much to invest in
each direction. Audit opportunities enter the action score of Stage~1, and
execution outcomes enter the next round of structural revision. The three
search decisions therefore share the same execution evidence and the same
posterior model, rather than being optimized separately and then combined.
This design lets limited evaluations revise search structure and capacity
jointly, avoiding wasted budget on low-potential directions and preventing
a newly admitted structure from being under-explored.

\begin{figure}[htbp]
\centering
\fbox{\parbox{0.93\linewidth}{\centering\color{red}
Framework placeholder: expert initialization $\rightarrow$ evidence-guided
search-state alternatives $\rightarrow$ audits, quotas, and joint revision
$\rightarrow$ candidate selection and execution $\rightarrow$ feedback.}}
\caption{Execution feedback connects search-state revision and allocation.}
\end{figure}

\begin{algorithm}[htbp]
\caption{JointTune: joint search-state revision}
\label{alg:JointTune}
\begin{algorithmic}[1]
\Require $\mathcal X,E,T,N_g,N_r,m,d_a,\gamma$, shared IID design $X_0$
\State $L\gets\textsc{Evaluate}(X_0)$; $\mathcal D\gets\textsc{Successful}(L)$
\State $(A,R,M)\gets\textsc{Initialize}(E,L,\mathcal X,d_a)$
\State $w_{\rm end}\gets\lceil\gamma(T-|X_0|)\rceil$
\For{$t=|X_0|,\ldots,T-1$}
 \State $P_t\gets\textsc{FitGP}(\mathcal D)$; $\bar x\gets\textsc{BestConfig}(\mathcal D)$
 \State $\mathcal V_t\gets\textsc{Configs}(L)$; $M\gets\textsc{ResetContext}(M,\bar x)$
 \State $\mathcal K\gets\{a_0\}$
 \If{$t<T-w_{\rm end}$}
  \State $\mathcal K\gets\mathcal K\cup\textsc{BestReshape}(A,R,P_t,M)$
  \If{$\textsc{ReviewDue}(t)$}
   \State $\mathcal K\gets\mathcal K\cup\textsc{BestExchangeAndCombined}(A,R,P_t,M)$
  \EndIf
 \EndIf
 \State $\mathcal P\gets\textsc{Experts}(E,L,\textsc{Evidence}(\mathcal K,P_t))$
 \State $\mathcal P\gets\textsc{ValidateOrFallback}(\mathcal P,R,\bar x)$
 \State $\mathcal Z\gets\textsc{AuditStates}(\mathcal K,\mathcal P,P_t,N_g,m,\mathcal V_t)$
 \State $(A',R',b,\mathcal A,m_t)\gets\mathcal Z[a_0]$
 \If{$t<T-w_{\rm end}\land\textsc{Calibrated}(M)\land m_t>0$}
  \State $(A',R',b,\mathcal A,m_t)\gets\textsc{JointSelect}(\mathcal Z,P_t)$
 \EndIf
 \State $(A,R)\gets(A',R')$; $M\gets\textsc{RecordState}(M,A,R,t)$
 \State $\mathcal Q_t\gets\textsc{Compile}(A,R,\bar x,\mathcal P)$
 \State $\mathcal C_g\gets\textsc{Fill}(\mathcal Q_t,b,\mathcal A,\mathcal V_t)$
 \State $\mathcal C_r\gets\textsc{Reserve}(N_r,\mathcal Q_t,M,\mathcal V_t)$
 \State $\mathcal C\gets\textsc{Deduplicate}((\mathcal C_g\cup\mathcal C_r)\setminus\mathcal V_t)$
 \State $x\gets\arg\max_{x'\in\mathcal C}a_t(x')$
 \If{$t<T-w_{\rm end}\land\textsc{Calibrated}(M)\land m_t>0$}
  \State $\mathcal S\gets\textsc{Shortlist}(\mathcal C)$
  \State $\bigl(\widehat{\mathcal I}_t(x),\sigma_{\rm MC}(x)\bigr)_{x\in\mathcal S}\gets\textsc{Info}(P_t,\mathcal A,\mathcal S)$
  \State $j\gets\arg\max_{x'\in\mathcal S}\widehat{\mathcal I}_t(x')$
  \If{$\widehat{\mathcal I}_t(j)>2\sigma_{\rm MC}(j)$} \State $x\gets j$ \EndIf
 \EndIf
 \If{$\textsc{ProbeDue}(t,M)\land t<T-w_{\rm end}$}
  \State $x\gets\textsc{ReservedProbe}(\mathcal C,M)$
 \EndIf
 \State $(y,s)\gets\textsc{Execute}(x)$; $L\gets L\cup\{(x,y,s)\}$
 \If{$s=\mathrm{SUCCESS}$}
  \State $\mathcal D\gets\mathcal D\cup\{(x,y)\}$; $M\gets\textsc{UpdateEvidence}(M,x,y,P_t,\mathcal A)$
 \Else
  \State $M.\mathrm{failed}\gets M.\mathrm{failed}\cup\{x\}$
 \EndIf
\EndFor
\State \Return $\textsc{BestConfig}(\mathcal D)$
\end{algorithmic}
\end{algorithm}

\subsection{Expert-Guided Initialization}\label{sec:initialization}

Four instances of the same LLM specialize in memory, concurrency and
execution, storage and WAL, and query optimization. They receive legal
domains, hardware/workload descriptions, and the default plus nine shared
IID measurements. Each returns regions, sparse patches, and
high/medium/low priority ratings on a common knob set.

MCTuner ranks knobs using weighted expert scores~\cite{mctuner}.
Relying only on such scores risks self-promotion: an expert may rate a knob
highly even when uncertain about it within its own group, which can suppress
knobs owned by other experts. JointTune therefore separates priority from
disagreement. Entropy provides an objective, cross-expert measure of overall
uncertainty for each knob, preventing highly uncertain knobs from entering
tuning prematurely. With $n_{j,\ell}$ valid ratings at level $\ell$, total
$n_j$, and smoothing $\alpha$, define
\begin{equation}
 \widehat p_{j,\ell}=\frac{n_{j,\ell}+\alpha}{n_j+3\alpha},\qquad
 H_j^{\rm sel}=-\frac{\sum_\ell\widehat p_{j,\ell}
                              \log\widehat p_{j,\ell}}{\log3}.
 \label{eq:ct-initial-entropy}
\end{equation}
Priority is $s_j=\widehat p_{j,\mathrm{high}}+
0.5\widehat p_{j,\mathrm{medium}}$. After protected knobs and domain
representatives, a quarter of remaining active slots, rounded up, goes to
high-entropy knobs with $s_j\geq0.5$; other slots follow priority.
Posterior gain breaks ties. Missing ratings are excluded, and unused
verification slots return to priority selection. Entropy preserves
opportunities to test disputed judgments; subsequent decisions depend on
execution evidence. Initial regions use bounded numerical centers/widths
or legal discrete preferences.

\subsection{Joint Search-State Revision}\label{sec:knob-selection}

We first specify the shared model and the budget accounting. A shared GP
models full configurations with a Mat\'ern-$5/2$ ARD kernel, normalized
numerical and ordinal coordinates, categorical indicators, and activation
masks. Experts see the complete execution history and propose legal changes;
the GP supplies the common utility evidence used to compare their proposals.
The execution budget $T$ counts workload evaluations, including initial
configurations, failed attempts, and reserved trials. Candidate capacity
$N=N_g+N_r$ separates group-mode quotas from reference and coverage channels;
each group quota includes $m$ audit slots, with $N_g\geq b_{\min}\geq m\geq1$.
Mode admission enforces
\begin{equation}
\begin{aligned}
 1\leq|\mathcal H_t|&\leq\left\lfloor N_g/b_{\min}\right\rfloor,\\
 N_{\rm rem}&=N_g-|\mathcal H_t|b_{\min}\geq0,\qquad
 \sum_hb_{h,t}=N_g.
\end{aligned}
\label{eq:ct-accounting}
\end{equation}
Modes with identical groups are merged. Posterior probes consume
computation, not execution slots.

Legal probes of group $C$ modify its knobs around $\bar x_t$, keeping
other coordinates fixed. For a joint posterior utility draw $U$, let
$\Delta_{C,t}=\max_x[U(x)-U(\bar x_t)]$ over those probes.
Conditional opportunity comprises $p_{C,t}=\Pr(\Delta_{C,t}>\delta)$ and
$g_{C,t}=\mathbb E[(\Delta_{C,t})_+]$, conditioned on $\mathcal D_t$.
Throughput uses utility $y=\log(f/f_0)$, and latency uses
$y=\log(f_0/f)$; $\delta=0.01$ is approximately a one-percent gain.

The action menu has four slots: keep; reshape a target incumbent; exchange
it for the strongest feasible inactive group; and combine the exchange
with reshaping a retained partner. On exchange reviews, the target is the
lowest-$g$ retirement-eligible group; otherwise it is the lowest-$g$ active
group. Alternatives are ordered by decreasing $g$. Within each slot,
retain the legal variant with highest preliminary equal-probe gain;
ties use fewer changes, then knob identifiers. Missing slots are omitted.
An exchange may admit at most a bounded number of knobs and must preserve
protected knobs, activation parents, and supported joint groups. It is
considered periodically; region-only actions remain available between
reviews.

For each action $a$, form equal-size mode audits
$\mathcal A_{h,t}(a)$ and compute its quota $b_h(a)$ by
Equation~\ref{eq:ct-budget-update}. Unchanged distributions reuse the same
audits across actions. First, measure the positive utility gain available
in each mode and aggregate it using the state's candidate shares:
\begin{equation}
\begin{aligned}
 D_h(a;U)&=\max_{x\in\mathcal A_{h,t}(a)}[U(x)-U(\bar x_t)-\delta]_+,\\
 G_t(a;U)&=\sum_h\frac{b_h(a)}{N_g}D_h(a;U).
\end{aligned}
\label{eq:ct-state-opportunity}
\end{equation}
$G_t$ is the expected best audited gain of a mode drawn according to
candidate shares. It changes with both the region and the capacity assigned
to that region. Next, compare each action with $a_0$, which retains the
current knobs and regions, using paired posterior draws:
\begin{equation}
\begin{aligned}
 \mu_t(a)&=\mathbb E[G_t(a;U)-G_t(a_0;U)\mid\mathcal D_t],\\
 \pi_t(a)&=\Pr(G_t(a;U)>G_t(a_0;U)\mid\mathcal D_t).
\end{aligned}
\label{eq:ct-action-advantage}
\end{equation}
Mean advantage measures the magnitude of the opportunity change;
superiority probability measures posterior support for its sign.
Retain actions passing both tests, then choose the strongest:
\begin{equation}
\begin{aligned}
 \mathcal K_t^+&=\{a_0\}\cup\{a\in\mathcal K_t:\mu_t(a)>\delta,\ \pi_t(a)\geq\rho\},\\
 a_t^*&=\arg\max_{a\in\mathcal K_t^+}\mu_t(a),\qquad\rho=0.8.
\end{aligned}
\label{eq:ct-joint-action}
\end{equation}
The chosen state is $(A(a_t^*),R(a_t^*),b(a_t^*))$. An exchange must also
beat the best admissible incumbent reshape by the same tests; weak results
in a restrictive range alone cannot justify retirement. Ties retain the
current state, then favor fewer changes.

An outgoing knob requires a minimum residence period, a minimum number of
successful local changes, distinct tested values, one fresh test after a
context reset, and no recent gain above $\delta$. Coverage and re-entry
trials are scheduled periodically; these executed trials supply evidence
for under-tested knobs, whereas generated or posterior-sampled points do
not. The final $\gamma$ fraction freezes structural revision and uses
improvement-based selection.

Coverage checks use predictions recorded before execution. For each
successful adaptive trial, store the standardized residual:
\begin{equation}
 e_i=\frac{|y_i-\mu_i^-|}{\sqrt{v_i^-+\tau_i^2}},\qquad
 \kappa_t=\frac{1}{|W_t|}\sum_{i\in W_t}\mathbf1\{e_i\leq1.96\}.
 \label{eq:ct-calibration}
\end{equation}
$\mu_i^-$ and $v_i^-$ are the latent predictive mean and variance,
$\tau_i^2$ is observation noise, and $W_t$ contains a bounded recent window
of successful predictions. $\textsc{Calibrated}(M_t)$ is true when at least
eight such predictions exist and $\kappa_t\geq0.7$. Otherwise we retain
$A_t,R_t$, use qLogNEI, and continue scheduled executed probes. Records
include gate pass rates, rejection reasons, and state-change frequencies.

\subsection{Evidence-Guided Regions and Proposals}
\label{sec:feedback-adaptation}

For numerical regions, legal probes cover an inner scale, the current scale
$r_j$, and an outer scale. Posterior draw $s$ has best probe $x^{*(s)}$.
Among draws with gain above $\delta$, let
$d_j^{(s)}=|z_j(x^{*(s)})-c_j|$, using normalized coordinate $z_j$ and
center $c_j$. Radius and conditional credible interval are
\begin{equation}
\begin{aligned}
 \widehat r_j&=\operatorname{clip}_{\mathcal R_j}Q_{0.8}(d_j^{(s)}),\\
 [r_j^-,r_j^+]&=\operatorname{clip}_{\mathcal R_j}
                  [Q_{0.1}(d_j^{(s)}),Q_{0.9}(d_j^{(s)})],\\
 \mathcal R_j&=[\max(r_{\min},r_{\rm in}r_j),\min(r_{\max},r_{\rm out}r_j)].
\end{aligned}
\label{eq:ct-radius}
\end{equation}
$Q_\tau$ denotes an empirical posterior quantile. If too few improving
draws are available, the current radius is retained and the mode is marked
for verification. The evidence packet gives radius change, interval, the
unconditional improvement probability, and paired action advantage.

Experts receive this packet in their regular request, approve actions with
mechanisms, and return at most a bounded number of action-tagged patches
each. Center shifts are bounded; edited centers are re-scored. Shrinking
also requires successful points on both sides of the center. Cross-domain
patches require agreement from their experts. Accepted patches change at
most a bounded number of actual knobs with bounded total normalized
movement, including feasibility repairs. Numerical sampling retains a small
global exploration fraction; discrete preferences enumerate legal values.

The controller compiles approved patches and regions into $\mathcal Q_t$,
which stores mode distributions, validated patch seeds, reviewed cross-domain
patches, and the selected scope/anchor. $\textsc{Fill}$ samples those
distributions; patches provide seeds within mode quotas and cause no further
LLM calls. An empty or failed expert response uses its last legal
distribution projected to the selected scope, or an anchor-centered legal
fallback. Every failure/fallback is recorded. Posterior references are
capped at a bounded number of distinct points, with shared audits reused
across actions. Optional combined and exchange variants are dropped in that
order if the cap is exceeded; an exchange is never compared without its
incumbent reshape. The comparison uses one fitted GP and joint posterior
draws, without additional mandatory LLM requests or workload executions.
Logs distinguish expansion, contraction, shift, hold, and each gate that
prevents adoption.

\subsection{Audited Allocation and Informative Trials}
\label{sec:candidate-selection}

Let $\mathcal V_t$ contain every previously attempted configuration in $L$,
including $M_t.\mathrm{failed}$. Auditing, $\textsc{Fill}$, and
$\textsc{Reserve}$ exclude $\mathcal V_t$ before admitting a candidate.
Each state generates up to $m$ legal, distinct audit points per mode.
Across actions and modes, use the minimum attainable audit count
$m_t$; if zero, keep the current state and use the reference channel.
No points are copied to fill a quota.
$\textsc{Reserve}(N_r,\mathcal Q_t,M_t,\mathcal V_t)$ supplies legal IID
reference points, reviewed cross-domain patches, and scheduled coverage or
re-entry points. It uses the compiled anchor/scope and memory of actual
coverage, residence, inactive waiting time, and failed attempts. Its union
with $\textsc{Fill}$ is deduplicated by canonical configuration key and
filtered against $\mathcal V_t$ again. Deduplication retains the source
set $S(x)$, and audits are frozen before scoring states or filling quotas.
For each posterior draw, let $x^\dagger$ be the best audit point. Set
$J_t=\varnothing$ if $U(x^\dagger)-U(\bar x_t)\leq\delta$, and
$J_t=S(x^\dagger)$ otherwise. In each state, compute
\begin{equation}
\begin{aligned}
 w_{h,t}&=\sum_{S:h\in S}\frac{\Pr(J_t=S\mid\mathcal D_t)}{|S|},\\
 b_{h,t}&=b_{\min}+\operatorname{RoundSum}_h\!\left[
 N_{\rm rem}\frac{\lambda+w_{h,t}}{\sum_g(\lambda+w_{g,t})}\right],
 \qquad\lambda=0.02.
\end{aligned}
\label{eq:ct-budget-update}
\end{equation}
Remainder-preserving rounding maintains total $N_g$. Audits already occupy
quota slots; additional generation does not enlarge them. Group credit is
split across contributing sources, and expert emphasis derives from
$v_e=\sum_{h:e\in E_h}b_h/|E_h|$. A changed distribution receives fresh
credit before its quota affects the structural decision.

On the selected state's candidate pool $\mathcal C_t$, let $a_t(x)$ be
noisy expected improvement, computed by qLogNEI with batch size one.
The shortlist contains a bounded number of points from
$\mathcal C_t^+=\{x:a_t(x)\geq\eta\max_{x'}a_t(x')\}$, $\eta=0.8$.
Within it, evaluate information about the frozen audit label:
\begin{equation}
 \mathcal I_t(x)=H(J_t\mid\mathcal D_t)-
 \mathbb E_{Y_x\mid\mathcal D_t}
       [H(J_t\mid\mathcal D_t,(x,Y_x))].
 \label{eq:ct-information}
\end{equation}
For $F$ independent noisy fantasies, let $R_f(x)$ be the prior label
entropy minus conditional label entropy, each estimated with latent draws
and common random numbers across candidates. The estimator and its outer
Monte Carlo standard error are
\begin{equation}
 \widehat{\mathcal I}_t(x)=\frac1F\sum_{f=1}^F R_f(x),\qquad
 \sigma_{\rm MC}(x)=\sqrt{\frac{\sum_f
 (R_f(x)-\widehat{\mathcal I}_t(x))^2}{F(F-1)}}.
 \label{eq:ct-mc-error}
\end{equation}
$\textsc{Info}$ returns one estimate/error pair
$(\widehat{\mathcal I}_t(x),\sigma_{\rm MC}(x))$ for every shortlisted $x$.
Choose the largest $\widehat{\mathcal I}_t$ only when calibrated, $m_t>0$,
and its estimate exceeds $2\sigma_{\rm MC}$; otherwise maximize qLogNEI.
This error estimates outer fantasy variability, separately from predictive
coverage. Better opportunity identification can redirect later candidate
shares; the EI guard retains strong immediate prospects while testing
allocation-relevant alternatives. Reserved coverage/re-entry trials have
priority, and the final $\gamma$ fraction uses qLogNEI.

Successful trials update observations, coverage, residuals, and audit
credit; failures remain in history without training the utility GP.
Algorithm~\ref{alg:JointTune} summarizes the loop.

% ============================================================
% TODO before submission (not part of the main text):
%
% 1. Replace the red framework placeholder with the final figure.
%
% 2. Add Table tab:ct-symbols defining X, E, T, n0, N, N_g, N_r, m,
%    b_min, N_rem, d_a, delta, eta, rho, lambda, alpha, gamma, kappa,
%    F, and the calibration window.
%
% 3. Add Table tab:ct-settings with all default values and controller
%    periods, including:
%      N=100, N_g=80, N_r=20, m=4, b_min=5;
%      mode cap and overflow retention rule;
%      gamma=0.1 and boundary handling of w_end;
%      review interval (every six rounds) and at most two new knobs;
%      coverage every third adaptive round, re-entry every eighth;
%      calibration window (last at most 20 predictions) and minimum 8;
%      radius inner/outer scales 0.8/1.5, bounds 0.025/0.5;
%      patch limit 4 knobs, movement limit 1.5, global exploration 10%;
%      posterior reference cap 128, shortlist 16, F=16 fantasies,
%      256 latent draws, fewer than 32 improving draws threshold.
%
% 4. Decide whether Algorithm 1 should show the mode-cap merge step
%    explicitly; currently it is described in the text only.
%
% 5. In the experiments, report the Calibrated gate pass rate, the
%    information-selection trigger rate, and a comparison against pure
%    qLogNEI. Also report LLM call counts/tokens/latency and GP fitting
%    and posterior-query overhead.
% ============================================================
\section{Experiment}\label{sec:experiment}

\subsection{Experiment Setup}\label{sec:setup}

\subsubsection{Target Workload}\label{sec:workload}

We evaluate JointTune on five representative benchmarks covering OLAP,
OLTP, and HTAP workloads. The OLAP workload is TPC-H, used to assess
latency optimization on complex analytical queries. The OLTP workloads
include TPC-C, Twitter, and YCSB, covering standard transaction
processing, high-concurrency social workloads, and read--write mixed
workloads, respectively; they are used to assess throughput improvement.
The HTAP workload is CH-benCHmark, used to evaluate both transactional
throughput and analytical query latency. These five benchmarks cover the
main database operation modes from analytical to transactional and from
single-purpose to hybrid, and are widely adopted in database tuning
studies. All implementations are sourced from the standard benchmark
library~\cite{difallah2013oltpbench,leis2015job}, ensuring consistency and reproducibility
of our setup.

\subsubsection{Hardware and Environment}\label{sec:hardware}

Experiments are conducted on a server equipped with two 12-core Intel
Xeon E5-2678 v3 processors running at 2.50 GHz, providing 24 physical
cores and 48 hardware threads, 128 GB of RAM, and a 2 TB hard disk drive
(HDD). The database runs in a Docker container with persistent
ext4-backed storage. LLM experts are accessed through the DeepSeek API.
All compared methods run on the same server, with the same database
version and the same container configuration, to avoid hardware and
environment differences.

\subsubsection{Baselines}\label{sec:baselines}

JointTune is implemented in Python, and its LLM experts are accessed
through the DeepSeek API. We compare against the following seven
representative tuning methods.

\begin{itemize}
  \item \textbf{SMAC}~\cite{smac}: a sequential model-based optimization
        method that adopts a random forest surrogate, which scales well to
        high-dimensional and categorical input.
  \item \textbf{GP}~\cite{ituned}: classical Gaussian-process Bayesian
        optimization, the underlying optimizer of iTuned, which links knob
        configurations with performance and updates the model with new
        observations.
  \item \textbf{CDBtune}~\cite{cdbtune}: an end-to-end cloud database
        tuning system using deep reinforcement learning, which employs DDPG
        to generate configurations in high-dimensional continuous space.
  \item \textbf{Restune}~\cite{restune}: a resource-oriented tuning method
        that casts configuration tuning as constrained Bayesian optimization
        and uses historical task data and meta-learning to accelerate cloud
        database tuning.
  \item \textbf{OtterTune}~\cite{ottertune}: a classic automatic tuning
        system based on workload characterization and Gaussian-process
        Bayesian optimization, which reuses historical tuning experience
        to recommend configurations.
  \item \textbf{MCTuner}~\cite{mctuner}: a tuning framework that combines
        LLM-based mixture-of-experts knob selection with MCTS spatial
        decomposition, representing the current LLM-guided line of work.
  \item \textbf{DOT}~\cite{dot}: a dynamic knob selection and online
        sampling algorithm that uses RFECV to prune low-importance knobs
        and a likelihood-ratio test to balance exploration and exploitation.
\end{itemize}

In summary, SMAC~\cite{smac}, GP~\cite{ituned}, Restune~\cite{restune},
and OtterTune~\cite{ottertune} are Bayesian-optimization based methods;
CDBtune~\cite{cdbtune} is a reinforcement-learning based method; and
MCTuner~\cite{mctuner} and DOT~\cite{dot} represent the LLM-guided and
dynamic-selection routes, respectively. All methods run under the same
execution budget. For LLM-based methods (MCTuner~\cite{mctuner} and
JointTune), expert-call overhead is recorded separately and is not
counted toward the workload execution budget; for GP-based methods, we
follow OtterTune's~\cite{ottertune} configuration-space processing and
surrogate-model settings.

\subsubsection{Metrics}\label{sec:metrics}

We adopt relative performance improvement as the primary evaluation
metric~\cite{e2etune}. For the OLAP workload TPC-H, the objective is to
minimize query latency, and the improvement is defined as
\begin{equation}
  \Delta = \frac{lat_{\mathrm{default}} - lat}{lat_{\mathrm{default}}},
  \label{eq:delta-latency}
\end{equation}
where $lat$ and $lat_{\mathrm{default}}$ denote the 95th-percentile
latency under the optimized and default configurations, respectively.
For the OLTP workloads TPC-C, Twitter, and YCSB, the objective is to
maximize throughput (transactions per second, $tps$), and the
improvement is defined as
\begin{equation}
  \Delta = \frac{tps - tps_{\mathrm{default}}}{tps_{\mathrm{default}}},
  \label{eq:delta-throughput}
\end{equation}
where $tps$ and $tps_{\mathrm{default}}$ denote the throughput under the
optimized and default configurations, respectively. For the HTAP
workload CH-benCHmark, we report both transactional throughput and the
95th-percentile latency of analytical queries to reflect performance
under mixed workloads. Lower latency and higher throughput are better;
all figures report relative improvement, where a larger value indicates
better performance.

\subsubsection{Tuning Settings}\label{sec:tuning-settings}

We use PostgreSQL v14.9 as the target DBMS, and LLM experts are accessed
through the DeepSeek API. Following prior work~\cite{gptuner,mctuner,cdbtune},
all methods start from the same 60 knobs and then apply their own
selection or range-compression mechanisms. For example, MCTuner and
DBTune use LLM knowledge for knob selection, OtterTune and Restune
initialize models from historical workloads or resource features, and
DOT updates the active knob set through dynamic pruning.

For a fair comparison under a controlled execution budget, each method is
limited to 100 real configuration evaluations, with one configuration
sampled and evaluated per iteration. For TPC-H, we optimize the
95th-percentile latency; for TPC-C, Twitter, and YCSB, we optimize
throughput; for CH-benCHmark, we optimize both transactional throughput
and the 95th-percentile latency of analytical queries. After each
evaluation, the target workload is executed, and the database is
restarted once a method proposes a new knob configuration, since some
parameters take effect only after a restart.

For Bayesian-optimization based methods (JointTune, SMAC, GP, OtterTune,
and Restune), the first 10 iterations use Latin hypercube sampling (LHS)
for cold-start initialization, and the remaining 90 iterations focus on
configuration optimization. For reinforcement-learning based methods
(e.g., CDBtune), we follow recent work and avoid retraining the neural
network on the target workload to reduce the risk of
overfitting~\cite{zhang2021facilitating}. For JointTune, the execution
budget $T$ counts only real configuration evaluations; candidate
generation, posterior probes, and LLM expert calls are not counted
toward $T$, but their call counts, token consumption, and latency are
recorded separately to compare the actual tuning overhead across methods.
